\documentclass[11pt, a4paper]{article}
\usepackage[a4paper, top=2.5cm, bottom=2.5cm, left=2cm, right=2cm]{geometry}
\usepackage{fontspec}
\usepackage[english, bidi=basic, provide=*]{babel}
\babelprovide[import, onchar=ids fonts]{english}
\babelfont{rm}{TeX Gyre Termes}
\babelfont{sf}{TeX Gyre Heros}
\babelfont{tt}{TeX Gyre Cursor}
\usepackage{enumitem}
\usepackage{amsmath}

\title{\textbf{Chapter 5: The Structure and Function of Large Biological Molecules \\ 75 Multiple Choice Questions}}
\author{}
\date{}

\begin{document}

\maketitle

\begin{enumerate}

% ---------------------------------------
% 5.1 Polymers and Monomers (1-7)
% ---------------------------------------

\item \textbf{Which of the following classes of large biological molecules does NOT consist of true polymers?}
\begin{enumerate}[label=\Alph*.]
    \item Proteins
    \item Carbohydrates
    \item Lipids
    \item Nucleic acids
    \item Polysaccharides
\end{enumerate}
\textbf{Answer:} C \\
\textbf{Explanation:} While lipids are large, complex biological molecules, they are not composed of repeating, uniform monomer units joined together in a chain like the other three classes.

\item \textbf{The chemical reaction that joins two monomers together by removing a molecule of water is known as:}
\begin{enumerate}[label=\Alph*.]
    \item Hydrolysis
    \item Hydration synthesis
    \item A dehydration reaction
    \item Oxidation
    \item Denaturation
\end{enumerate}
\textbf{Answer:} C \\
\textbf{Explanation:} In a dehydration reaction, two molecules are covalently bonded with the loss of a water molecule ($H_2O$).

\item \textbf{The process of digesting polymers into their constituent monomers is called:}
\begin{enumerate}[label=\Alph*.]
    \item Dehydration synthesis
    \item Hydrolysis
    \item Esterification
    \item Denaturation
    \item Polymerization
\end{enumerate}
\textbf{Answer:} B \\
\textbf{Explanation:} Hydrolysis ("water breakage") involves adding a water molecule to break the covalent bond between monomers.

\item \textbf{In a typical dehydration reaction used to synthesize a polymer, what two components are contributed by the reacting monomers?}
\begin{enumerate}[label=\Alph*.]
    \item A hydroxyl group and a hydrogen atom
    \item Two hydrogen atoms and one oxygen atom
    \item Two hydroxyl groups
    \item A carboxyl group and an amino group
    \item A phosphate group and a methyl group
\end{enumerate}
\textbf{Answer:} A \\
\textbf{Explanation:} One monomer contributes a hydroxyl group ($-OH$) and the other contributes a hydrogen ($-H$), yielding the released water molecule.

\item \textbf{Specialized macromolecules that speed up chemical reactions in cells, including the making and breaking of polymers, are called:}
\begin{enumerate}[label=\Alph*.]
    \item Hormones
    \item Enzymes
    \item Steroids
    \item Chaperonins
    \item Isomers
\end{enumerate}
\textbf{Answer:} B \\
\textbf{Explanation:} Enzymes are biological catalysts that facilitate processes like dehydration synthesis and hydrolysis without being consumed by the reaction.

\item \textbf{If a polymer consists of 100 monomers, how many molecules of water must be removed to link them together completely in a linear chain?}
\begin{enumerate}[label=\Alph*.]
    \item 100
    \item 99
    \item 101
    \item 50
    \item 200
\end{enumerate}
\textbf{Answer:} B \\
\textbf{Explanation:} Linking monomers into a linear polymer requires $n-1$ dehydration reactions, where $n$ is the total number of monomers.

\item \textbf{Which of the following best describes the structural foundation of molecular diversity in polymers?}
\begin{enumerate}[label=\Alph*.]
    \item A limited set of monomers can be arranged in nearly limitless, varied sequences.
    \item Cells continuously discover entirely new elements.
    \item All polymers are identical in sequence but differ in size.
    \item Monomers randomly mutate into new atomic structures.
    \item Polymers are entirely synthesized without enzymes.
\end{enumerate}
\textbf{Answer:} A \\
\textbf{Explanation:} Like the alphabet, a relatively small set of basic units (e.g., 20 amino acids) can be ordered in an infinite variety of ways to build uniquely functioning macromolecules.

% ---------------------------------------
% 5.2 Carbohydrates (8-22)
% ---------------------------------------

\item \textbf{The molecular formula of a typical monosaccharide generally represents a multiple of:}
\begin{enumerate}[label=\Alph*.]
    \item $CH_4$
    \item $CH_2O$
    \item $C_6H_{12}O_6$
    \item $NH_3$
    \item $H_2O$
\end{enumerate}
\textbf{Answer:} B \\
\textbf{Explanation:} Carbohydrates are literally "carbon water." A common example is glucose ($C_6H_{12}O_6$), which is exactly 6 times $CH_2O$.

\item \textbf{Depending on the location of its carbonyl group, a simple sugar can be classified as either a(n):}
\begin{enumerate}[label=\Alph*.]
    \item Hexose or a pentose
    \item Monomer or a polymer
    \item Aldose or a ketose
    \item Enantiomer or a diastereomer
    \item Steroid or a fat
\end{enumerate}
\textbf{Answer:} C \\
\textbf{Explanation:} If the carbonyl group is at the end of the skeleton, it is an aldehyde (aldose). If it is within the skeleton, it is a ketone (ketose).

\item \textbf{In aqueous solutions, how do five- and six-carbon sugars like glucose naturally arrange themselves?}
\begin{enumerate}[label=\Alph*.]
    \item They form straight, unbranched chains.
    \item They disintegrate into carbon dioxide and water.
    \item They form rigid ring structures.
    \item They instantly polymerize into starch.
    \item They bond tightly to dissolved nitrogen.
\end{enumerate}
\textbf{Answer:} C \\
\textbf{Explanation:} In water, the most stable, favored form for most 5- and 6-carbon sugars is a closed ring structure.

\item \textbf{A covalent bond formed between two monosaccharides by a dehydration reaction is known specifically as a(n):}
\begin{enumerate}[label=\Alph*.]
    \item Phosphodiester linkage
    \item Peptide bond
    \item Ester linkage
    \item Glycosidic linkage
    \item Sulfhydryl cross-link
\end{enumerate}
\textbf{Answer:} D \\
\textbf{Explanation:} A glycosidic linkage is the specific type of covalent bond that connects carbohydrate monomers.

\item \textbf{The disaccharide sucrose, ordinary table sugar, is composed of which two monosaccharides?}
\begin{enumerate}[label=\Alph*.]
    \item Glucose and Galactose
    \item Glucose and Glucose
    \item Fructose and Galactose
    \item Glucose and Fructose
    \item Maltose and Glucose
\end{enumerate}
\textbf{Answer:} D \\
\textbf{Explanation:} Sucrose is transported widely in plants and is formed by joining a glucose monomer to a fructose monomer.

\item \textbf{Plants store surplus glucose for later energy use in the form of:}
\begin{enumerate}[label=\Alph*.]
    \item Cellulose
    \item Glycogen
    \item Chitin
    \item Starch
    \item Cholesterol
\end{enumerate}
\textbf{Answer:} D \\
\textbf{Explanation:} Starch is a polymer of glucose monomers and represents stored energy. Plants stockpile starch as granules within cellular structures known as plastids.

\item \textbf{Animals store excess glucose primarily in liver and muscle cells as a highly branched polymer known as:}
\begin{enumerate}[label=\Alph*.]
    \item Starch
    \item Glycogen
    \item Amylopectin
    \item Cellulose
    \item Glucagon
\end{enumerate}
\textbf{Answer:} B \\
\textbf{Explanation:} Glycogen is extensively branched, allowing for rapid hydrolysis to release glucose into the bloodstream when energy is demanded.

\item \textbf{The major structural carbohydrate found in the tough walls of plant cells is:}
\begin{enumerate}[label=\Alph*.]
    \item Starch
    \item Glycogen
    \item Chitin
    \item Peptidoglycan
    \item Cellulose
\end{enumerate}
\textbf{Answer:} E \\
\textbf{Explanation:} Cellulose is the most abundant organic compound on Earth, providing critical structural support to plant cell walls.

\item \textbf{Why can humans digest starch but not cellulose?}
\begin{enumerate}[label=\Alph*.]
    \item Humans lack the enzymes that can hydrolyze the $\beta$ (beta) linkages of cellulose.
    \item Starch has fewer calories than cellulose.
    \item Cellulose is a lipid, not a carbohydrate.
    \item Cellulose is poisoned by stomach acid.
    \item Cellulose molecules form irreversible triple bonds.
\end{enumerate}
\textbf{Answer:} A \\
\textbf{Explanation:} Starch relies on $\alpha$ linkages, which human enzymes can easily break. Cellulose uses alternating $\beta$ linkages, forming straight molecules that human enzymes cannot recognize or hydrolyze.

\item \textbf{Which property is characteristic of cellulose molecules?}
\begin{enumerate}[label=\Alph*.]
    \item They are highly branched.
    \item They form helical, spring-like structures.
    \item They are straight and never branch, allowing them to hydrogen-bond with parallel cellulose molecules.
    \item They dissolve easily in water.
    \item They contain internal peptide bonds.
\end{enumerate}
\textbf{Answer:} C \\
\textbf{Explanation:} The $\beta$ linkages in cellulose force it into a straight shape. The hydroxyl groups on parallel molecules bond together, forming incredibly strong microfibrils.

\item \textbf{When a cow digests grass, how is the cellulose broken down?}
\begin{enumerate}[label=\Alph*.]
    \item Cows naturally produce specialized cellulose-digesting enzymes.
    \item Cows chew the grass until the mechanical force breaks the covalent bonds.
    \item Symbiotic bacteria and protists in the cow's rumen produce enzymes to hydrolyze the cellulose.
    \item Stomach acid instantly turns cellulose into glucose gas.
    \item Grass does not contain cellulose.
\end{enumerate}
\textbf{Answer:} C \\
\textbf{Explanation:} Like termites, cows harbor a rich microbial community in their digestive tracts that possesses the correct enzymes to break down cellulose into usable sugar.

\item \textbf{Which carbohydrate is used by arthropods to build their exoskeletons?}
\begin{enumerate}[label=\Alph*.]
    \item Cellulose
    \item Peptidoglycan
    \item Starch
    \item Chitin
    \item Glycogen
\end{enumerate}
\textbf{Answer:} D \\
\textbf{Explanation:} Chitin is a structural polysaccharide similar to cellulose, except its glucose monomers contain a nitrogen-containing appendage.

\item \textbf{In addition to arthropod exoskeletons, chitin is also found in:}
\begin{enumerate}[label=\Alph*.]
    \item The cell walls of fungi.
    \item The roots of woody plants.
    \item Human fingernails.
    \item Bacterial cell membranes.
    \item Plant seeds.
\end{enumerate}
\textbf{Answer:} A \\
\textbf{Explanation:} Fungi uniquely use chitin rather than cellulose as the primary structural building block of their cell walls.

\item \textbf{Lactose intolerance in humans is caused by:}
\begin{enumerate}[label=\Alph*.]
    \item An allergy to the proteins in milk.
    \item The inability to synthesize enough lactase, the enzyme that breaks the glycosidic linkage in lactose.
    \item The toxic effects of pure galactose.
    \item The rapid conversion of milk sugars into alcohol.
    \item A lack of bacteria in the stomach.
\end{enumerate}
\textbf{Answer:} B \\
\textbf{Explanation:} Without the enzyme lactase, lactose cannot be hydrolyzed. It moves into the large intestine, where bacteria ferment it, causing discomfort.

\item \textbf{A molecule consisting of 5 glucose units joined by glycosidic linkages would require the removal of how many water molecules during its synthesis?}
\begin{enumerate}[label=\Alph*.]
    \item 1
    \item 4
    \item 5
    \item 6
    \item 10
\end{enumerate}
\textbf{Answer:} B \\
\textbf{Explanation:} Joining $n$ monomers requires $n-1$ dehydration reactions. Joining 5 units requires 4 linkages, thus removing 4 water molecules.

% ---------------------------------------
% 5.3 Lipids (23-38)
% ---------------------------------------

\item \textbf{Which characteristic defines all lipids?}
\begin{enumerate}[label=\Alph*.]
    \item They are all comprised of complex ring structures.
    \item They mix poorly, if at all, with water.
    \item They all contain nitrogen.
    \item They are all exactly identical in size.
    \item They are highly hydrophilic.
\end{enumerate}
\textbf{Answer:} B \\
\textbf{Explanation:} Lipids are universally grouped by their hydrophobic nature, stemming from their long, nonpolar hydrocarbon regions.

\item \textbf{A typical fat (triglyceride) is composed of:}
\begin{enumerate}[label=\Alph*.]
    \item Three glycerols and one fatty acid.
    \item One glycerol and three fatty acids.
    \item A phosphate group and two fatty acids.
    \item Four fused carbon rings.
    \item Amino acids and steroids.
\end{enumerate}
\textbf{Answer:} B \\
\textbf{Explanation:} In making a fat, three long-chain fatty acids are covalently joined to a single three-carbon glycerol molecule.

\item \textbf{The covalent bond connecting a fatty acid to a glycerol molecule in a fat is called a(n):}
\begin{enumerate}[label=\Alph*.]
    \item Glycosidic linkage
    \item Peptide bond
    \item Phosphodiester bond
    \item Ester linkage
    \item Disulfide bridge
\end{enumerate}
\textbf{Answer:} D \\
\textbf{Explanation:} An ester linkage is formed by a dehydration reaction between a hydroxyl group (on glycerol) and a carboxyl group (on the fatty acid).

\item \textbf{Why are fats highly hydrophobic?}
\begin{enumerate}[label=\Alph*.]
    \item Because they contain many polar hydroxyl groups.
    \item Because water molecules hydrogen-bond to one another and exclude the relatively nonpolar hydrocarbon chains.
    \item Because the ester linkages repel oxygen.
    \item Because they carry a strong negative charge.
    \item Because they are exclusively solid.
\end{enumerate}
\textbf{Answer:} B \\
\textbf{Explanation:} The nonpolar C-H bonds in the long fatty acid tails cannot form hydrogen bonds, so water forces the fat molecules together to minimize disruption of water's hydrogen bond network.

\item \textbf{A fatty acid with no double bonds between carbon atoms in its hydrocarbon chain is said to be:}
\begin{enumerate}[label=\Alph*.]
    \item Polyunsaturated
    \item Monounsaturated
    \item Saturated
    \item Trans
    \item Amphipathic
\end{enumerate}
\textbf{Answer:} C \\
\textbf{Explanation:} It is "saturated" with the maximum possible number of hydrogen atoms, because all carbon-carbon bonds are single bonds.

\item \textbf{Most animal fats (e.g., lard, butter) have saturated fatty acid tails, which makes them:}
\begin{enumerate}[label=\Alph*.]
    \item Gases at room temperature.
    \item Solids at room temperature because the straight tails pack closely together.
    \item Liquids at room temperature because they easily slide past each other.
    \item Highly soluble in water.
    \item Blue in color.
\end{enumerate}
\textbf{Answer:} B \\
\textbf{Explanation:} The absence of double bonds means the chains are straight, allowing the triglyceride molecules to pack tightly into a solid matrix.

\item \textbf{What structural feature creates a "kink" or bend in an unsaturated fatty acid?}
\begin{enumerate}[label=\Alph*.]
    \item A \textit{cis} double bond between two carbon atoms.
    \item An extra methyl group on the tail.
    \item The addition of a phosphate group.
    \item A \textit{trans} triple bond.
    \item An inserted oxygen atom.
\end{enumerate}
\textbf{Answer:} A \\
\textbf{Explanation:} Nearly every double bond in naturally occurring fatty acids is a \textit{cis} double bond, which creates a rigid angle (kink) in the hydrocarbon chain.

\item \textbf{Unsaturated fats like olive oil are liquid at room temperature because:}
\begin{enumerate}[label=\Alph*.]
    \item They have exceptionally short carbon chains.
    \item The kinks in their \textit{cis} double bonds prevent tight packing.
    \item They are heavily hydrated by water.
    \item They lack ester linkages.
    \item They are completely non-polar.
\end{enumerate}
\textbf{Answer:} B \\
\textbf{Explanation:} The structural kinks physically prevent the molecules from packing closely enough to solidify at room temperature.

\item \textbf{The process of "hydrogenation" (creating hydrogenated vegetable oils) results in:}
\begin{enumerate}[label=\Alph*.]
    \item Converting saturated fats to unsaturated fats.
    \item Breaking down fats into separate fatty acids and glycerol.
    \item Synthetically adding hydrogen to unsaturated fats to convert them into saturated fats (and creating \textit{trans} fats).
    \item Producing synthetic proteins.
    \item Creating highly unstable, explosive liquids.
\end{enumerate}
\textbf{Answer:} C \\
\textbf{Explanation:} Hydrogenation forces hydrogen into unsaturated oils, straightening the chains to make them solid at room temperature (like peanut butter or margarine) but also creating harmful \textit{trans} double bonds.

\item \textbf{Dietary \textit{trans} fats have been heavily linked to which major health issue?}
\begin{enumerate}[label=\Alph*.]
    \item Cystic fibrosis
    \item Cardiovascular disease (atherosclerosis)
    \item Scurvy
    \item Malaria
    \item Osteoporosis
\end{enumerate}
\textbf{Answer:} B \\
\textbf{Explanation:} Trans fats contribute to atherosclerosis (plaque buildup in blood vessels) even more strongly than naturally saturated fats.

\item \textbf{What is the primary biological function of fats?}
\begin{enumerate}[label=\Alph*.]
    \item Catalyzing chemical reactions
    \item Carrying genetic information
    \item Short-term energy storage
    \item Long-term energy storage
    \item Defending against viruses
\end{enumerate}
\textbf{Answer:} D \\
\textbf{Explanation:} A gram of fat stores more than twice as much energy as a gram of a polysaccharide like starch, making it an excellent, compact, long-term energy reservoir.

\item \textbf{Phospholipids are unique among lipids because they are "amphipathic," meaning:}
\begin{enumerate}[label=\Alph*.]
    \item They have two distinct heads.
    \item They contain both a hydrophilic region and a hydrophobic region.
    \item They exist in a state between liquid and gas.
    \item They are composed of both proteins and lipids.
    \item They act as both an acid and a base.
\end{enumerate}
\textbf{Answer:} B \\
\textbf{Explanation:} The phosphate-containing head is highly polar (hydrophilic), while the two hydrocarbon tails are nonpolar (hydrophobic).

\item \textbf{When phospholipids are added to water, they spontaneously self-assemble into a double-layered sheet called a "bilayer." Why?}
\begin{enumerate}[label=\Alph*.]
    \item The hydrophobic tails attract water, while the heads repel it.
    \item The hydrophobic tails point inward away from water, while the hydrophilic heads point outward facing the water.
    \item They form strong covalent bonds with each other.
    \item The water evaporates, leaving the lipids behind.
    \item Enzymes in the water construct the bilayer.
\end{enumerate}
\textbf{Answer:} B \\
\textbf{Explanation:} This structural arrangement minimizes the energy of the system by shielding the hydrophobic tails entirely from the aqueous environment, creating the basis of all cell membranes.

\item \textbf{Steroids are lipids characterized by a carbon skeleton consisting of:}
\begin{enumerate}[label=\Alph*.]
    \item Three fused rings.
    \item Four fused rings.
    \item Long straight chains entirely lacking rings.
    \item A helical shape.
    \item Two attached phosphate groups.
\end{enumerate}
\textbf{Answer:} B \\
\textbf{Explanation:} All steroids, despite diverse functions, share this core structure of four fused carbon rings.

\item \textbf{Which of the following is an example of a crucial biological steroid found in animal cell membranes?}
\begin{enumerate}[label=\Alph*.]
    \item Triglyceride
    \item Estrogen
    \item Cellulose
    \item Cholesterol
    \item Phospholipid
\end{enumerate}
\textbf{Answer:} D \\
\textbf{Explanation:} Cholesterol is both a vital component of animal cell membranes and the precursor from which vertebrate sex hormones are synthesized.

\item \textbf{Although cholesterol is essential for life, a high level of it in the blood is correlated with:}
\begin{enumerate}[label=\Alph*.]
    \item Low blood pressure
    \item Immune deficiency
    \item Atherosclerosis
    \item Rapid muscle growth
    \item Increased bone density
\end{enumerate}
\textbf{Answer:} C \\
\textbf{Explanation:} High levels of cholesterol, often exacerbated by saturated and trans fat consumption, contribute to the arterial plaques of cardiovascular disease.

% ---------------------------------------
% 5.4 Proteins (39-58)
% ---------------------------------------

\item \textbf{Proteins account for more than \_\_\_\_\_ of the dry mass of most cells.}
\begin{enumerate}[label=\Alph*.]
    \item 10\%
    \item 25\%
    \item 50\%
    \item 75\%
    \item 90\%
\end{enumerate}
\textbf{Answer:} C \\
\textbf{Explanation:} Proteins are incredibly ubiquitous and are involved in almost every dynamic function of living beings.

\item \textbf{Enzymes are a specific class of proteins that function to:}
\begin{enumerate}[label=\Alph*.]
    \item Store genetic information.
    \item Form the tough fibers of hair and nails.
    \item Selectively accelerate chemical reactions without being consumed.
    \item Transport oxygen in the blood.
    \item Provide insulation to the body.
\end{enumerate}
\textbf{Answer:} C \\
\textbf{Explanation:} Enzymatic proteins act as catalysts, repeatedly driving essential cellular reactions forward.

\item \textbf{The monomers that are joined together to construct a protein polymer are:}
\begin{enumerate}[label=\Alph*.]
    \item Nucleotides
    \item Monosaccharides
    \item Fatty acids
    \item Amino acids
    \item Isomers
\end{enumerate}
\textbf{Answer:} D \\
\textbf{Explanation:} A protein is a biologically functional molecule made of one or more polypeptides, which are in turn polymers of amino acids.

\item \textbf{At the center of an amino acid is an asymmetric carbon atom called the $\alpha$ (alpha) carbon. What four partners are attached to it?}
\begin{enumerate}[label=\Alph*.]
    \item An amino group, a carboxyl group, a hydrogen atom, and a variable R group.
    \item An amino group, a phosphate group, a hydrogen atom, and a methyl group.
    \item Two hydroxyl groups, a hydrogen atom, and an R group.
    \item Four identical variable R groups.
    \item A nucleotide base, a sugar, a phosphate, and a hydrogen.
\end{enumerate}
\textbf{Answer:} A \\
\textbf{Explanation:} This standard core structure is shared by all 20 amino acids; only the variable side chain (R group) differs.

\item \textbf{The physical and chemical properties of the \_\_\_\_\_\_\_ determine the unique characteristics of a particular amino acid.}
\begin{enumerate}[label=\Alph*.]
    \item Alpha carbon
    \item Amino group
    \item Carboxyl group
    \item R group (side chain)
    \item Hydrogen atom
\end{enumerate}
\textbf{Answer:} D \\
\textbf{Explanation:} The R group can be nonpolar, polar, acidic, or basic, dictating how the amino acid will behave and fold in a polypeptide.

\item \textbf{If an amino acid's R group consists solely of carbon and hydrogen atoms, the amino acid is most likely:}
\begin{enumerate}[label=\Alph*.]
    \item Hydrophilic
    \item Highly acidic
    \item Hydrophobic (nonpolar)
    \item Charged
    \item A base
\end{enumerate}
\textbf{Answer:} C \\
\textbf{Explanation:} Hydrocarbons are nonpolar. Amino acids with nonpolar side chains are hydrophobic and tend to cluster in the interior of a folded protein.

\item \textbf{Amino acids with acidic side chains are generally negatively charged at cellular pH. These acidic side chains primarily feature which functional group?}
\begin{enumerate}[label=\Alph*.]
    \item Methyl
    \item Sulfhydryl
    \item Carboxyl
    \item Hydroxyl
    \item Amino
\end{enumerate}
\textbf{Answer:} C \\
\textbf{Explanation:} The carboxyl group donates an $H^+$ to the solution, leaving the side chain with a net negative charge.

\item \textbf{The covalent bond connecting the carboxyl group of one amino acid to the amino group of the next is called a:}
\begin{enumerate}[label=\Alph*.]
    \item Glycosidic linkage
    \item Peptide bond
    \item Hydrogen bond
    \item Phosphodiester linkage
    \item Ester bond
\end{enumerate}
\textbf{Answer:} B \\
\textbf{Explanation:} Peptide bonds form through a dehydration reaction, creating a repeating polypeptide backbone with extending R groups.

\item \textbf{The two distinct ends of a polypeptide chain are known as the:}
\begin{enumerate}[label=\Alph*.]
    \item 3' end and 5' end
    \item Alpha end and Beta end
    \item N-terminus (amino end) and C-terminus (carboxyl end)
    \item Hydrophilic end and Hydrophobic end
    \item Cis and Trans ends
\end{enumerate}
\textbf{Answer:} C \\
\textbf{Explanation:} One end has a free amino group (N-terminus), and the opposite end has a free carboxyl group (C-terminus).

\item \textbf{The specific, linear sequence of amino acids in a protein constitutes its:}
\begin{enumerate}[label=\Alph*.]
    \item Primary structure
    \item Secondary structure
    \item Tertiary structure
    \item Quaternary structure
    \item Quintessential structure
\end{enumerate}
\textbf{Answer:} A \\
\textbf{Explanation:} Primary structure is simply the exact order of amino acids, which is directly dictated by genetic information.

\item \textbf{A slight change in a protein's primary structure, such as the single amino acid substitution that causes sickle-cell disease, can:}
\begin{enumerate}[label=\Alph*.]
    \item Have absolutely no effect on protein function.
    \item Cause the protein to turn into a lipid.
    \item Drastically affect its 3D shape and its ability to function normally.
    \item Convert the protein into a nucleic acid.
    \item Double the speed of enzymatic reactions.
\end{enumerate}
\textbf{Answer:} C \\
\textbf{Explanation:} Substituting valine for glutamic acid alters hemoglobin's shape drastically, demonstrating how primary structure dictates higher levels of structure and ultimate function.

\item \textbf{Secondary structures, like the $\alpha$ helix and $\beta$ pleated sheet, are the result of \_\_\_\_\_\_\_\_\_ between the repeating constituents of the polypeptide backbone.}
\begin{enumerate}[label=\Alph*.]
    \item Ionic bonds
    \item Disulfide bridges
    \item Hydrogen bonds
    \item Covalent bonds
    \item Hydrophobic exclusions
\end{enumerate}
\textbf{Answer:} C \\
\textbf{Explanation:} Oxygen atoms in the backbone have a partial negative charge, and hydrogens attached to nitrogens have a partial positive charge, allowing regular hydrogen bonds to form along the backbone.

\item \textbf{Which secondary structure is a delicate coil held together by hydrogen bonding between every fourth amino acid?}
\begin{enumerate}[label=\Alph*.]
    \item $\beta$ pleated sheet
    \item Double helix
    \item $\alpha$ helix
    \item Tertiary tangle
    \item Chaperonin
\end{enumerate}
\textbf{Answer:} C \\
\textbf{Explanation:} The alpha helix is a common secondary structure found heavily in fibrous proteins like keratin (hair).

\item \textbf{Tertiary structure—the overall 3D shape of a polypeptide—is primarily determined by interactions between:}
\begin{enumerate}[label=\Alph*.]
    \item The atoms of the repeating backbone.
    \item The R groups (side chains) of the various amino acids.
    \item Two distinct polypeptide chains.
    \item RNA and the ribosome.
    \item The solvent and the glass container.
\end{enumerate}
\textbf{Answer:} B \\
\textbf{Explanation:} While secondary structure involves the backbone, tertiary structure is driven by hydrophobic interactions, ionic bonds, hydrogen bonds, and covalent links between the side chains.

\item \textbf{As a polypeptide folds into its functional shape in water, amino acids with nonpolar side chains typically cluster at the core of the protein. This is known as a(n):}
\begin{enumerate}[label=\Alph*.]
    \item Disulfide bridge
    \item Ionic bond
    \item Hydrophobic interaction
    \item Enzymatic reaction
    \item Hydration shell
\end{enumerate}
\textbf{Answer:} C \\
\textbf{Explanation:} Water molecules hydrogen-bond to each other and to hydrophilic parts of the protein, essentially "excluding" nonpolar side chains and forcing them into the center.

\item \textbf{Strong covalent bonds that reinforce the protein's overall shape, forming between two cysteine monomers, are called:}
\begin{enumerate}[label=\Alph*.]
    \item Hydrogen bonds
    \item Peptide bonds
    \item Glycosidic linkages
    \item Disulfide bridges
    \item Ester linkages
\end{enumerate}
\textbf{Answer:} D \\
\textbf{Explanation:} The sulfhydryl groups ($-SH$) on two cysteines react to form a durable $-S-S-$ bond, riveting parts of the protein together.

\item \textbf{Quaternary structure occurs ONLY when a protein consists of:}
\begin{enumerate}[label=\Alph*.]
    \item A single polypeptide chain folded in on itself.
    \item Two or more individual polypeptide chains aggregated into one functional macromolecule.
    \item Both DNA and RNA.
    \item At least 1,000 amino acids.
    \item Both a lipid and a carbohydrate.
\end{enumerate}
\textbf{Answer:} B \\
\textbf{Explanation:} Proteins like collagen (3 intertwined polypeptides) and hemoglobin (4 subunits) have quaternary structure resulting from multiple distinct polypeptide chains coming together.

\item \textbf{When a protein loses its native shape due to changes in pH, salt concentration, or excessive heat, the process is called:}
\begin{enumerate}[label=\Alph*.]
    \item Renaturation
    \item Hydrolysis
    \item Denaturation
    \item Catalysis
    \item Dehydration
\end{enumerate}
\textbf{Answer:} C \\
\textbf{Explanation:} Environmental changes disrupt the weak chemical bonds holding the protein's shape, causing it to unravel and become biologically inactive.

\item \textbf{Why does an extremely high fever present a dangerous health risk?}
\begin{enumerate}[label=\Alph*.]
    \item It causes DNA to instantly mutate.
    \item Excessive heat can denature critical cellular proteins and enzymes.
    \item It turns saturated fats into trans fats.
    \item It hydrolyzes cellulose in the body.
    \item It creates spontaneous disulfide bridges everywhere.
\end{enumerate}
\textbf{Answer:} B \\
\textbf{Explanation:} As temperatures rise dangerously high, thermal agitation disrupts hydrogen bonds and other weak interactions, unraveling proteins.

\item \textbf{Proteins whose primary function is to assist in the proper folding of other newly synthesized proteins are called:}
\begin{enumerate}[label=\Alph*.]
    \item Antibodies
    \item Receptors
    \item Chaperonins (chaperone proteins)
    \item Lipoproteins
    \item Ligases
\end{enumerate}
\textbf{Answer:} C \\
\textbf{Explanation:} Chaperonins do not specify the final structure; rather, they provide an isolated, perfect environment for the new polypeptide to fold spontaneously without disruptive chemical interference from the cytoplasm.

% ---------------------------------------
% 5.5 Nucleic Acids (59-71)
% ---------------------------------------

\item \textbf{The primary function of nucleic acids is to:}
\begin{enumerate}[label=\Alph*.]
    \item Provide immediate energy to cells.
    \item Form the structural barriers of cell membranes.
    \item Store, transmit, and help express hereditary information.
    \item Transport oxygen.
    \item Catalyze digestive reactions.
\end{enumerate}
\textbf{Answer:} C \\
\textbf{Explanation:} DNA and RNA carry the "blueprints" for assembling proteins, dictating cellular functions.

\item \textbf{The correct flow of genetic information in a typical cell (Gene Expression) is:}
\begin{enumerate}[label=\Alph*.]
    \item $Protein \rightarrow RNA \rightarrow DNA$
    \item $RNA \rightarrow DNA \rightarrow Protein$
    \item $DNA \rightarrow Protein \rightarrow RNA$
    \item $DNA \rightarrow RNA \rightarrow Protein$
    \item $Protein \rightarrow DNA \rightarrow RNA$
\end{enumerate}
\textbf{Answer:} D \\
\textbf{Explanation:} DNA provides instructions to synthesize messenger RNA (mRNA), which then interacts with cellular machinery (ribosomes) to direct the assembly of proteins.

\item \textbf{The monomer building block of a nucleic acid is a:}
\begin{enumerate}[label=\Alph*.]
    \item Nucleoside
    \item Amino acid
    \item Monosaccharide
    \item Nucleotide
    \item Phospholipid
\end{enumerate}
\textbf{Answer:} D \\
\textbf{Explanation:} A nucleotide is the repeating unit of a polynucleotide strand.

\item \textbf{What are the three components of every nucleotide?}
\begin{enumerate}[label=\Alph*.]
    \item A sugar, an amino group, and a carboxyl group.
    \item A six-carbon sugar, a purine, and a pyrimidine.
    \item A pentose sugar, a nitrogenous base, and one to three phosphate groups.
    \item A nitrogenous base, a glycerol, and a fatty acid.
    \item A steroid ring, a phosphate group, and a hydroxyl group.
\end{enumerate}
\textbf{Answer:} C \\
\textbf{Explanation:} Nucleotides contain a 5-carbon sugar (pentose), a nitrogen-containing base, and at least one phosphate group.

\item \textbf{Nitrogenous bases are divided into pyrimidines and purines. Which of the following is characteristic of a purine?}
\begin{enumerate}[label=\Alph*.]
    \item It has a single six-membered ring.
    \item It consists of a six-membered ring fused to a five-membered ring.
    \item It lacks nitrogen completely.
    \item It is only found in RNA.
    \item It is identical to a pentose sugar.
\end{enumerate}
\textbf{Answer:} B \\
\textbf{Explanation:} Purines (Adenine and Guanine) are larger, double-ring structures, whereas pyrimidines (Cytosine, Thymine, Uracil) are smaller single-ring structures.

\item \textbf{Which specific nitrogenous base is found ONLY in DNA, and never in RNA?}
\begin{enumerate}[label=\Alph*.]
    \item Adenine (A)
    \item Guanine (G)
    \item Cytosine (C)
    \item Thymine (T)
    \item Uracil (U)
\end{enumerate}
\textbf{Answer:} D \\
\textbf{Explanation:} In RNA, Thymine is replaced by Uracil. The other three (A, G, C) are shared by both.

\item \textbf{What is the difference between the sugars found in DNA (deoxyribose) and RNA (ribose)?}
\begin{enumerate}[label=\Alph*.]
    \item Deoxyribose has six carbons; ribose has five.
    \item Deoxyribose is missing an oxygen atom on the second carbon ring compared to ribose.
    \item Ribose is missing an entire phosphate group.
    \item Deoxyribose is linear; ribose is a ring.
    \item Ribose has a nitrogenous base attached; deoxyribose does not.
\end{enumerate}
\textbf{Answer:} B \\
\textbf{Explanation:} The prefix "deoxy-" literally means it lacks an oxygen atom that is present on the 2' carbon of ribose.

\item \textbf{Adjacent nucleotides in a polynucleotide are joined by:}
\begin{enumerate}[label=\Alph*.]
    \item Peptide bonds
    \item Glycosidic linkages
    \item Phosphodiester linkages
    \item Ester linkages
    \item Hydrogen bonds entirely
\end{enumerate}
\textbf{Answer:} C \\
\textbf{Explanation:} A phosphodiester linkage consists of a phosphate group linking the sugars of two nucleotides, creating a repeating sugar-phosphate backbone.

\item \textbf{The two ends of a polynucleotide strand have built-in directionality. What chemical group is located at the 5' end?}
\begin{enumerate}[label=\Alph*.]
    \item A free hydroxyl group
    \item A free amino group
    \item A free phosphate group
    \item A free carboxyl group
    \item A bare nitrogenous base
\end{enumerate}
\textbf{Answer:} C \\
\textbf{Explanation:} The 5' end has a phosphate group attached to the 5' carbon of the sugar. The opposite 3' end has a free hydroxyl group attached to the 3' carbon.

\item \textbf{In a DNA double helix, the two sugar-phosphate backbones run in opposite 5' $\rightarrow$ 3' directions from each other. This arrangement is termed:}
\begin{enumerate}[label=\Alph*.]
    \item Antiparallel
    \item Parallel
    \item Perpendicular
    \item Semi-conservative
    \item Symmetrical
\end{enumerate}
\textbf{Answer:} A \\
\textbf{Explanation:} Like a divided highway, the structural strands run in opposite directions to allow stable base pairing across the center.

\item \textbf{According to the strict rules of complementary base pairing in DNA, Adenine (A) always pairs with:}
\begin{enumerate}[label=\Alph*.]
    \item Guanine (G)
    \item Cytosine (C)
    \item Thymine (T)
    \item Uracil (U)
    \item Another Adenine (A)
\end{enumerate}
\textbf{Answer:} C \\
\textbf{Explanation:} Adenine (a purine) always pairs with Thymine (a pyrimidine) via two hydrogen bonds.

\item \textbf{Because of complementary base pairing, if a DNA strand has the sequence 5'-AGCT-3', the opposite strand will be:}
\begin{enumerate}[label=\Alph*.]
    \item 3'-TCGA-5'
    \item 5'-TCGA-3'
    \item 3'-AGCT-5'
    \item 5'-CTAG-3'
    \item 3'-UCGU-5'
\end{enumerate}
\textbf{Answer:} A \\
\textbf{Explanation:} A pairs with T, G pairs with C. The strands are antiparallel, so the matching sequence runs 3' to 5'.

\item \textbf{Which statement regarding RNA is typically true?}
\begin{enumerate}[label=\Alph*.]
    \item It exists strictly as a double helix.
    \item It utilizes Thymine instead of Uracil.
    \item It usually exists as a single polynucleotide strand.
    \item It contains the sugar deoxyribose.
    \item It is only found within the nucleus.
\end{enumerate}
\textbf{Answer:} C \\
\textbf{Explanation:} Unlike the rigid double helix of DNA, RNA usually exists as a single strand, though it can fold onto itself (e.g., tRNA).

% ---------------------------------------
% 5.6 Genomics and Proteomics (72-75)
% ---------------------------------------

\item \textbf{Analyzing large sets of genes or comparing whole genomes of different species is an approach called:}
\begin{enumerate}[label=\Alph*.]
    \item Proteomics
    \item Cytology
    \item Bioinformatics
    \item Genomics
    \item Taxonomy
\end{enumerate}
\textbf{Answer:} D \\
\textbf{Explanation:} Genomics involves the comprehensive study of an organism's entire genome and the comparison of genomes between species.

\item \textbf{The vast amount of sequencing data generated by modern biology requires the use of computer software and computational tools. This field is called:}
\begin{enumerate}[label=\Alph*.]
    \item Cybernetics
    \item Proteomics
    \item Bioinformatics
    \item Astrobiology
    \item Synthetic biology
\end{enumerate}
\textbf{Answer:} C \\
\textbf{Explanation:} Bioinformatics blends biology, computer science, and information technology to process huge biological datasets.

\item \textbf{By analyzing the genomes of various species, scientists have found that two species that appear to be very closely related evolutionarily (like humans and macaques) will:}
\begin{enumerate}[label=\Alph*.]
    \item Have completely unrelated DNA sequences.
    \item Share a greater proportion of their DNA and protein sequences than more distantly related species.
    \item Have identical genomes.
    \item Produce completely different amino acids.
    \item Not share any common ancestors.
\end{enumerate}
\textbf{Answer:} B \\
\textbf{Explanation:} Molecular genealogy confirms anatomical and fossil evidence: close evolutionary relatives share high sequence similarity due to descent from a recent common ancestor.

\item \textbf{Proteomics focuses on:}
\begin{enumerate}[label=\Alph*.]
    \item The structure of protons in biology.
    \item Identifying the order of base pairs in DNA.
    \item Analyzing large sets of proteins, including their sequences and interactions.
    \item The synthesis of lipids.
    \item Generating synthetic ribosomes.
\end{enumerate}
\textbf{Answer:} C \\
\textbf{Explanation:} Just as genomics studies large-scale gene data, proteomics analyzes the full sets of proteins expressed by cells to understand their functions and interactions.

\end{enumerate}

\end{document}
